% GENERATED FILE. Edit canonical lessons, then run npm run generate:books.
% canonical-source-hash: fnv1a64:b57498dfceafac0c
% canonical-lessons: GU-C154-kapdu, GU-C154-sadi, GU-C154-dupatto, GU-C154-utarvu, GU-C154-odhvu

\chapter{Saris, scarves and cloth}
\label{ch:gu-saris-scarves-and-cloth}
\hlchaptermodality{154}

\begin{chapteropening}
\textbf{By the end of this chapter:} \emph{I can talk about the clothes people put on and take off.}

Complete the last lesson of chapter 154: I can talk about the clothes people put on and take off.
\end{chapteropening}

\section[\texorpdfstring{\gu{કપડું} (kapḍũ)}{કપડું (kapḍũ)}]{\gu{કપડું} (kapḍũ) — cloth, garment}

\label{lesson:GU-C154-kapdu}



\begin{quote}
Before the new one: say the Gujarati for veranda, then the Gujarati for stairs.
\end{quote}

\subsection*{You'll want to know: \gu{કપડું}}
\textbf{\gu{કપડું}} — \emph{kapḍũ} — \textquotedblleft{}cloth, garment\textquotedblright{}.

\begin{quote}
\textbf{\gu{આ} \gu{કપડું} \gu{નરમ} \gu{છે}.} — \emph{ā kapḍũ naram che} — \textquotedblleft{}This cloth is soft.\textquotedblright{}
\end{quote}

The plural \textbf{\gu{કપડાં}}, \emph{kapḍā̃}, means clothes.

\begin{cousinweb}
\textbf{\gu{કપડું}} (\emph{kapḍũ}) is a tadbhava, a word inherited from Sanskrit and reshaped by sound change. It comes from \emph{karpaṭa}, \textquotedblleft{}patched garment, rag\textquotedblright{}; CDIAL glosses the Gujarati word \textquotedblleft{}any garment\textquotedblright{} and asks whether it came by way of Hindi (CDIAL 2871).
\end{cousinweb}

\begin{grammarlens}[title={What you've built}]
One more word for everyday life at A2.
\end{grammarlens}

\subsection*{Your turn}
\begin{itemize}
\raggedright
  \item \emph{Say it:} \emph{kapḍũ}
  \item \emph{Say it:} \emph{kapḍũ}, once more
  \item \emph{Say it:} \emph{ā kapḍũ naram che} — This cloth is soft.
  \item \emph{Recall:} say the Gujarati for veranda, then the Gujarati for stairs, then say \emph{kapḍũ} again
\end{itemize}

\begin{tcolorbox}[breakable,colback=teal!4,colframe=teal!35!black,title={Before you move on}]
What does \gu{કપડું} mean? (\textquotedblleft{}Cloth, garment\textquotedblright{}.) Say it once more.
\end{tcolorbox}
\section[\texorpdfstring{\gu{સાડી} (sāḍī)}{સાડી (sāḍī)}]{\gu{સાડી} (sāḍī) — sari}

\label{lesson:GU-C154-sadi}



\begin{quote}
Before the new one: say the Gujarati for stairs, then the Gujarati for cloth.
\end{quote}

\subsection*{You'll want to know: \gu{સાડી}}
\textbf{\gu{સાડી}} — \emph{sāḍī} — \textquotedblleft{}sari\textquotedblright{}.

\begin{quote}
\textbf{\gu{મારી} \gu{બહેન} \gu{સાડી} \gu{પહેરે} \gu{છે}.} — \emph{mārī bahen sāḍī pahere che} — \textquotedblleft{}My sister wears a sari.\textquotedblright{}
\end{quote}

\begin{cousinweb}
\textbf{\gu{સાડી}} (\emph{sāḍī}) is a tadbhava. It comes from Sanskrit \emph{śāṭa}, \emph{śāṭī}, \textquotedblleft{}strip of cloth\textquotedblright{}, and CDIAL glosses the Gujarati word \textquotedblleft{}sari\textquotedblright{} (CDIAL 12381).
\end{cousinweb}

\begin{grammarlens}[title={What you've built}]
One more word for everyday life at A2.
\end{grammarlens}

\subsection*{Your turn}
\begin{itemize}
\raggedright
  \item \emph{Say it:} \emph{sāḍī}
  \item \emph{Say it:} \emph{sāḍī}, once more
  \item \emph{Say it:} \emph{mārī bahen sāḍī pahere che} — My sister wears a sari.
  \item \emph{Recall:} say the Gujarati for stairs, then the Gujarati for cloth, then say \emph{sāḍī} again
\end{itemize}

\begin{tcolorbox}[breakable,colback=teal!4,colframe=teal!35!black,title={Before you move on}]
What does \gu{સાડી} mean? (\textquotedblleft{}Sari\textquotedblright{}.) Say it once more.
\end{tcolorbox}
\section[\texorpdfstring{\gu{દુપટ્ટો} (dupaṭṭo)}{દુપટ્ટો (dupaṭṭo)}]{\gu{દુપટ્ટો} (dupaṭṭo) — scarf, stole}

\label{lesson:GU-C154-dupatto}



\begin{quote}
Before the new one: say the Gujarati for cloth, then the Gujarati for sari.
\end{quote}

\subsection*{You'll want to know: \gu{દુપટ્ટો}}
\textbf{\gu{દુપટ્ટો}} — \emph{dupaṭṭo} — \textquotedblleft{}scarf, stole\textquotedblright{}.

\begin{quote}
\textbf{\gu{તેનો} \gu{દુપટ્ટો} \gu{લાલ} \gu{છે}.} — \emph{teno dupaṭṭo lāl che} — \textquotedblleft{}Her scarf is red.\textquotedblright{}
\end{quote}

\begin{cousinweb}
\textbf{\gu{દુપટ્ટો}} (\emph{dupaṭṭo}) is built from \emph{du-}, \textquotedblleft{}two\textquotedblright{}, and \emph{paṭṭo}, \textquotedblleft{}strip of cloth\textquotedblright{}. Under Sanskrit \emph{paṭṭa}, \textquotedblleft{}cloth\textquotedblright{}, which it thinks probably non-Aryan, CDIAL lists the kindred Gujarati \emph{pāṭ}, \textquotedblleft{}strip of cloth\textquotedblright{}, and \emph{pāṭo}, \textquotedblleft{}bandage\textquotedblright{}, but no form of the compound itself (CDIAL 7700).
\end{cousinweb}

\begin{grammarlens}[title={What you've built}]
One more word for everyday life at A2.
\end{grammarlens}

\subsection*{Your turn}
\begin{itemize}
\raggedright
  \item \emph{Say it:} \emph{dupaṭṭo}
  \item \emph{Say it:} \emph{dupaṭṭo}, once more
  \item \emph{Say it:} \emph{teno dupaṭṭo lāl che} — Her scarf is red.
  \item \emph{Recall:} say the Gujarati for cloth, then the Gujarati for sari, then say \emph{dupaṭṭo} again
\end{itemize}

\begin{tcolorbox}[breakable,colback=teal!4,colframe=teal!35!black,title={Before you move on}]
What does \gu{દુપટ્ટો} mean? (\textquotedblleft{}Scarf, stole\textquotedblright{}.) Say it once more.
\end{tcolorbox}
\section[\texorpdfstring{\gu{ઉતારવું} (utārvũ)}{ઉતારવું (utārvũ)}]{\gu{ઉતારવું} (utārvũ) — to take off, to take down}

\label{lesson:GU-C154-utarvu}



\begin{quote}
Before the new one: say the Gujarati for sari, then the Gujarati for scarf.
\end{quote}

\subsection*{You'll want to know: \gu{ઉતારવું}}
\textbf{\gu{ઉતારવું}} — \emph{utārvũ} — \textquotedblleft{}to take off, to take down\textquotedblright{}.

\begin{quote}
\textbf{\gu{તે} \gu{ઘરે} \gu{બૂટ} \gu{ઉતારે} \gu{છે}.} — \emph{te ghare būṭ utāre che} — \textquotedblleft{}He takes off his shoes at home.\textquotedblright{}
\end{quote}

\begin{cousinweb}
\textbf{\gu{ઉતારવું}} (\emph{utārvũ}) is a tadbhava. It comes from Sanskrit \emph{uttārayati}, \textquotedblleft{}makes come out, makes descend\textquotedblright{}, and CDIAL glosses the Gujarati verb \textquotedblleft{}to set down, bring down\textquotedblright{} (CDIAL 1770).
\end{cousinweb}

\begin{grammarlens}[title={What you've built}]
One more word for everyday life at A2.
\end{grammarlens}

\subsection*{Your turn}
\begin{itemize}
\raggedright
  \item \emph{Say it:} \emph{utārvũ}
  \item \emph{Say it:} \emph{utārvũ}, once more
  \item \emph{Say it:} \emph{te ghare būṭ utāre che} — He takes off his shoes at home.
  \item \emph{Recall:} say the Gujarati for sari, then the Gujarati for scarf, then say \emph{utārvũ} again
\end{itemize}

\begin{tcolorbox}[breakable,colback=teal!4,colframe=teal!35!black,title={Before you move on}]
What does \gu{ઉતારવું} mean? (\textquotedblleft{}To take off, to take down\textquotedblright{}.) Say it once more.
\end{tcolorbox}
\section[\texorpdfstring{\gu{ઓઢવું} (oḍhvũ)}{ઓઢવું (oḍhvũ)}]{\gu{ઓઢવું} (oḍhvũ) — to wrap oneself in, to cover oneself}

\label{lesson:GU-C154-odhvu}



\begin{quote}
Before the new one: say the Gujarati for scarf, then the Gujarati for to take off.
\end{quote}

\subsection*{You'll want to know: \gu{ઓઢવું}}
\textbf{\gu{ઓઢવું}} — \emph{oḍhvũ} — \textquotedblleft{}to wrap oneself in, to cover oneself\textquotedblright{}.

\begin{quote}
\textbf{\gu{શિયાળામાં} \gu{હું} \gu{ધાબળો} \gu{ઓઢું} \gu{છું}.} — \emph{śiyāḷāmā̃ hũ dhābaḷo oḍhũ chũ} — \textquotedblleft{}In winter I wrap myself in a blanket.\textquotedblright{}
\end{quote}

\begin{cousinweb}
\textbf{\gu{ઓઢવું}} (\emph{oḍhvũ}) is inherited through Prakrit. It comes from \emph{*ōḍḍh}, \textquotedblleft{}put on, wear, cover\textquotedblright{}, a reconstructed form that CDIAL takes to be borrowed from Dravidian, beside Tamil \emph{uṭu-}, \textquotedblleft{}to dress\textquotedblright{}, though it notes that Burrow derives the Tamil verb otherwise (CDIAL 2547).
\end{cousinweb}

\begin{grammarlens}[title={What you've built}]
One more word for everyday life at A2.
\end{grammarlens}

\subsection*{Your turn}
\begin{itemize}
\raggedright
  \item \emph{Say it:} \emph{oḍhvũ}
  \item \emph{Say it:} \emph{oḍhvũ}, once more
  \item \emph{Say it:} \emph{śiyāḷāmā̃ hũ dhābaḷo oḍhũ chũ} — In winter I wrap myself in a blanket.
  \item \emph{Recall:} say the Gujarati for scarf, then the Gujarati for to take off, then say \emph{oḍhvũ} again
\end{itemize}

\begin{tcolorbox}[breakable,colback=teal!4,colframe=teal!35!black,title={Before you move on}]
What does \gu{ઓઢવું} mean? (\textquotedblleft{}To wrap oneself in, to cover oneself\textquotedblright{}.) Say it once more.
\end{tcolorbox}
